\documentclass[11pt,a4paper]{article}

\usepackage[utf8]{inputenc}
\usepackage[T2A]{fontenc}
\usepackage[russian,english]{babel}
\usepackage{amsmath,amssymb,amsfonts}
\usepackage{hyperref}
\usepackage{geometry}
\geometry{margin=2.5cm}

\title{GRA Repos Systematizer:\\
A Mathematical Index of the GRA / Hybrid Resonance / Multiverse Stack}

\author{qqewq}
\date{\today}

\begin{document}

\maketitle

\selectlanguage{english}

\begin{abstract}
We propose a mathematical view on the \emph{GRA Repos Systematizer}, 
a central index for approximately eighty GitHub repositories that form the
GRA / Hybrid Resonance / Multiverse ecosystem.
We formalize the repository set as a structured object, introduce
indexing and categorization operators, and relate them to the underlying
GRA axis: state space, subjectivity operators, meta-nulling, multiverse
trajectories, resonance and swarm-level dynamics.
This text is intended both as conceptual documentation and as a
lightweight reference for future formal work.
\end{abstract}

\selectlanguage{russian}

\begin{abstract}
В работе предлагается математический взгляд на \emph{GRA Repos Systematizer}~---
центральный индекс примерно восьмидесяти GitHub-репозиториев, образующих
экосистему GRA / Hybrid Resonance / Multiverse.
Множество репозиториев рассматривается как структурированный объект; вводятся
операторы индексирования и категоризации и связываются с осью GRA:
пространство состояний, операторы субъективности, мета-обнуление,
мультивселенские траектории, резонанс и роевой уровень динамики.
Текст одновременно служит концептуальной документацией и лёгкой
формальной ссылкой для дальнейших исследований.
\end{abstract}

\selectlanguage{english}

\section{Introduction}

The GRA / Hybrid Resonance / Multiverse stack is implemented as a
distributed family of repositories maintained under the
\texttt{github.com/qqewq} account.
Individually, these repositories explore subjectivity layers, meta-nulling
operators, multiverse alignment, swarm intelligence, physical embodiments
and culturally grounded constructs.
Taken together, they form a conceptual universe that is difficult to
grasp from source code alone.

The \emph{GRA Repos Systematizer} is a dedicated repository whose mission
is to provide a structured, mathematically informed index of this universe.
Rather than treating repositories as a flat list of URLs, we model them as
elements of a typed set equipped with categorization, projection and
aggregation operators.
This paper introduces a compact formalization of that index.

\section{Repository universe as a structured set}

\subsection{The set of repositories}

Let $\mathcal{R}$ denote the set of all repositories that belong to the
GRA ecosystem under consideration:
\[
  \mathcal{R} = \{ r_1, r_2, \dots, r_N \},
\]
where each $r_i$ is uniquely identified by its GitHub URL,
e.g.\ \texttt{https://github.com/qqewq/gra-core},
\texttt{https://github.com/qqewq/GRA-Subjectivity-Layer}, etc.
In the current state of the systematizer we have $N \approx 80$.

Each repository $r \in \mathcal{R}$ carries several attributes:
name, URL, informal description, primary theme
(e.g.\ subjectivity, meta-nulling, multiverse, swarm, ethnos),
and its position in the overall GRA architecture
(core, experimental, application, sandbox).
We denote the set of possible categories by $\mathcal{C}$.

\subsection{Index and category maps}

The central object of the systematizer is the \emph{index map}
\[
  I : \mathcal{R} \rightarrow \mathcal{C},
\]
which assigns to each repository $r$ a category $c = I(r) \in \mathcal{C}$.
In practice, $I$ is realized in the \texttt{README.md} of the
\texttt{gra-repos-systematizer} repository as a curated, human-readable
table of sections and lists.

We also consider a richer labeling
\[
  L : \mathcal{R} \rightarrow \mathcal{C} \times \mathcal{M}_{\text{axis}},
\]
where $\mathcal{M}_{\text{axis}}$ encodes alignment with elements of
the GRA conceptual axis (state space, subjectivity, meta-nulling,
multiverse, resonance).
Formally, for a repository $r$ we can write
\[
  L(r) = \bigl( c(r), a(r) \bigr),
\]
where $c(r) \in \mathcal{C}$ is a coarse category and
$a(r) \in \mathcal{M}_{\text{axis}}$ is a vector of booleans or weights
indicating how strongly $r$ contributes to different dimensions of the axis.

\section{GRA axis: states, subjectivity, multiverse, resonance}

\subsection{State space and trajectories}

We start from a state space $\mathcal{S}$ which encompasses cognitive,
environmental and cultural states relevant to GRA agents.
A trajectory is a finite sequence
\[
  \tau = (s_0, s_1, \dots, s_T), \quad s_t \in \mathcal{S}.
\]
The multiverse is then modeled as the set of admissible trajectories
\[
  \mathcal{M} = \{ \tau_i \mid \tau_i \text{ is an admissible trajectory} \}.
\]

Repositories containing ``Multiverse'' and ``Alignment'' in their names
implement various algorithms and experimental setups for generating,
evaluating and aligning trajectories $\tau \in \mathcal{M}$ with respect
to target functionals and constraints.

\subsection{Subjectivity as an operator}

Let $\mathcal{X}$ be the space of raw information (observations, logs,
external signals).
We define a subjectivity operator
\[
  \mathcal{O}_{\text{subj}} : \mathcal{X} \rightarrow \mathcal{S}_{\text{subj}},
\]
where $\mathcal{S}_{\text{subj}} \subseteq \mathcal{S}$ is the subspace
of subjectively relevant states.
Intuitively, $\mathcal{O}_{\text{subj}}$ implements the agent-specific
filtering, compression and recoding of the world into a subjective
configuration.

Repositories such as \texttt{GRA-Subjectivity-Layer} and
\texttt{GRA-Subjective-Evolution} explore how these operators can be
constructed, parameterized and evolved over time.

\subsection{Meta-nulling and zeroing}

Let $\Theta$ denote the space of subjective parameters:
beliefs, weights, priorities and other internal degrees of freedom.
We model meta-nulling as an operator
\[
  \mathcal{O}_{\text{null}} : \Theta \rightarrow \Theta,
\]
for which there exists a set of reference (``null'') configurations
$\Theta_0 \subseteq \Theta$ such that
\[
  \mathcal{O}_{\text{null}}(\theta) \rightarrow \theta_0 \in \Theta_0
\]
under iterated application.
Informally, meta-nulling aims to detect contradictions and traps and to
move the system towards a more neutral, meta-stable attractor.

Repositories \texttt{gra-meta-nullification},
\texttt{GRA-Meta-Nulling-Foundations},
\texttt{GRA-Meta-Nulling-Multiverse},
\texttt{GRA-Meta-Zeroing-Chaos},
\texttt{GRA-Paradox-Zeroing},
\texttt{gra-zeroing}
constitute different realizations and experiments around
$\mathcal{O}_{\text{null}}$.

\subsection{Resonance and swarm dynamics}

The set of agents $\mathcal{A} = \{a_1,\dots,a_N\}$ is treated as a
resonant swarm, where each agent has internal dynamics $F_i$ and a set
of interaction functions $G_{ij}$ with others.
At the level of the global state we can write
\[
  s_{t+1} = F(s_t) + R(s_t),
\]
where $F$ is the intrinsic evolution and $R$ is the resonance component
that aggregates cross-agent influences.

Swarm-oriented repositories (\texttt{GRA-Swarm},
\texttt{gra-swarm-agi}, \texttt{GRA-Swarm-Arena},
\texttt{GRA-Swarm-Engine}, \texttt{gra\_swarm\_asi},
\texttt{GRA\_DroneSwarm}) explore different forms of $R$ and relate them
to emergent intelligence and physical embodiments.

\section{Systematizer as a meta-layer}

\subsection{Index as a meta-operator}

On top of the GRA axis we place the systematizer as a meta-layer:
it does not introduce new dynamics on $\mathcal{S}$ directly, but
organizes the code-level artifacts implementing that dynamics.
Formally, we can view the index map $I$ as an operator
\[
  \mathcal{O}_{\text{index}} : \mathcal{R} \rightarrow \mathcal{C},
\]
with the following properties:
\begin{itemize}
  \item stability under small changes (minor updates of a repository
        should not change its category);
  \item extensibility (new elements $r_{N+1}$ can be assigned categories
        without breaking the existing structure);
  \item alignment with the conceptual axis (for each $c \in \mathcal{C}$
        there exists an interpretation in terms of
        $\mathcal{S}, \mathcal{M}, \mathcal{O}_{\text{subj}},
         \mathcal{O}_{\text{null}}, R$).
\end{itemize}

\subsection{Bilingual documentation}

The \texttt{gra-repos-systematizer} repository is designed to be
bilingual (Russian/English), both in its README and in supporting
documents such as this paper.
This reflects the cultural and linguistic context of GRA and allows the
same mathematical core to be communicated across different communities.

\section{Conclusion}

We have outlined a compact mathematical description of the
\emph{GRA Repos Systematizer} as a structured index over the GRA /
Hybrid Resonance / Multiverse repository universe.
The aim is not to fully axiomatize the stack, but to give a clear axis
along which individual repositories can be placed and related.
Future work may refine the label space $\mathcal{M}_{\text{axis}}$,
introduce richer morphisms between repositories and interpret certain
transitions as functors in a categorical setting.

\end{document}